\documentclass{beamer}
\usetheme{Madrid}
\usecolortheme{whale}

\title[FlixPulse Recommendation Dashboard]{FlixPulse: Movie Recommendation System \& Dashboard}
\subtitle{Cult Open Projects 2026 --- AI/ML Section}
\author[Cultural Council]{Cultural Council, IIT Roorkee}
\date{\today}

\begin{document}

% Slide 1: Title
\begin{frame}
  \titlepage
\end{frame}

% Slide 2: Problem Overview
\begin{frame}{1. Problem Overview \& Motivation}
  \begin{itemize}
    \item \textbf{Objective}: Design a personalized content discovery engine to boost subscriber retention.
    \item \textbf{Data Sparsity}: Matrix sparsity makes finding matching correlations difficult.
    \item \textbf{Rating Bias}: Skewed review distributions make models optimistic/pessimistic.
    \item \textbf{Evaluation Gap}: RMSE evaluates absolute errors, but users value ranking quality.
    \item \textbf{Explainability}: Transparency is needed to increase subscriber trust.
  end style
\end{frame}

% Slide 3: EDA & Dataset Characteristics
\begin{frame}{2. Dataset Features \& EDA Highlights}
  \begin{block}{Dense Subset Properties}
    \begin{itemize}
      \item \textbf{Dimension}: 21,045 active users, 1,287 movies, 5.7M ratings.
      \item \textbf{Matrix Sparsity}: 78.96\% (Density = 21.04\%).
    \end{itemize}
  \end{block}
  \begin{block}{Key Insights}
    \begin{itemize}
      \item \textbf{Positive Selection Skew}: Mean rating is 3.73 (4★ and 5★ reviews cover 63\% of total).
      \item \textbf{User Behavior}: Activity follows a heavy-tailed profile (most users rated 200--300 movies).
      \item \textbf{Long-Tail Blockbusters}: Popular titles capture a massive share of review counts.
    \end{itemize}
  \end{block}
\end{frame}

% Slide 4: Model Architectures
\begin{frame}{3. Model Design: SVD vs. Item-Based KNN}
  \begin{block}{A. Singular Value Decomposition (SVD)}
    \begin{itemize}
      \item Latent factor model: maps users and items to $f=50$ hidden features.
      \item Prediction: $\hat{r}_{u,i} = \mu + b_u + b_i + p_u^T q_i$.
      \item Optimized via Stochastic Gradient Descent (SGD) with L2 regularization.
    \end{itemize}
  \end{block}
  \begin{block}{B. Item-Based Collaborative Filtering (KNN)}
    \begin{itemize}
      \item Neighborhood model: calculates similarity weights between items.
      \item Similarity Metric: Adjusted Cosine Correlation.
      \item Prediction: Weighted average of user's ratings for adjacent items.
    \end{itemize}
  \end{block}
\end{frame}

% Slide 5: Experimental Results
\begin{frame}{4. Model Evaluation \& Metric Comparison}
  \begin{table}
    \centering
    \begin{tabular}{lcc}
      \toprule
      \textbf{Metric} & \textbf{SVD Model} & \textbf{Item-Based KNN} \\
      \midrule
      \textbf{Test RMSE} (Lower is better) & \textbf{0.8160} & 0.8553 \\
      \textbf{Test MAE} (Lower is better) & \textbf{0.6368} & 0.6687 \\
      \textbf{MAP@10} (Higher is better) & \textbf{0.7427} & 0.7079 \\
      \textbf{NDCG@10} (Higher is better) & \textbf{0.8337} & 0.8081 \\
      \textbf{Hit Rate@10} & \textbf{99.94\%} & 99.92\% \\
      \textbf{Diversity (ILD)} & \textbf{0.0477} & 0.0452 \\
      \textbf{Training Time} & \textbf{8.92s} & 225.86s \\
      \bottomrule
    \end{tabular}
  \end{table}
  \begin{itemize}
    \item \textbf{Accuracy \& Ranking}: SVD outperforms KNN across all vectors due to latent smoothing.
    \item \textbf{Inference Speed}: SVD predictions update in 0.31s vs. KNN's 4.85s.
  end style
\end{frame}

% Slide 6: Explainable recommendations
\begin{frame}{5. Explainable Recommendations Framework}
  \begin{block}{The Explainability Trade-off}
    SVD is highly accurate but uninterpretable. KNN computes explicit item similarity weights.
  \end{block}
  \begin{block}{Our Hybrid Solution}
    Use SVD to compute top-ranked recommendations, and KNN similarity matrices to generate natural explanations.
  \end{block}
  \begin{block}{Sample Output}
    \textit{``We recommend 'Toy Story 2' because you gave 'Toy Story' a high rating of 5 (Similarity: 98\%)''}
  \end{block}
\end{frame}

% Slide 7: Cold Start Strategy
\begin{frame}{6. Cold Start Strategy}
  \begin{block}{New Users}
    \begin{itemize}
      \item No history available for collaborative filters.
      \item \textbf{Fallback}: Popularity-based recommendations (highest average rating among movies with $\ge 2,000$ ratings).
    \end{itemize}
  \end{block}
  \begin{block}{New Movies}
    \begin{itemize}
      \item No ratings available for factorization.
      \item \textbf{Fallback 1}: Match initial similarities using genre overlaps.
      \item \textbf{Fallback 2}: Implement $\epsilon$-greedy exploration to randomly seed new uploads into active feeds.
    \end{itemize}
  \end{block}
\end{frame}

% Slide 8: Deployment & Future Scope
\begin{frame}{7. Decoupled Web Application \& Dashboard}
  \begin{itemize}
    \item \textbf{Backend API}: FastAPI (Python) serving uvicorn routes for predictions, similarities, search, and dashboard statistics.
    \item \textbf{Frontend Interface}: React SPA styled with dark-themed HSL glassmorphism design tokens.
    \item \textbf{Interactive Visuals}: Dynamic SVG charts plotting ratings distribution and model performance.
    \item \textbf{Future Scope}: Integrate deep Neural Collaborative networks and TMDB scraper APIs for visual movie poster assets.
  end style
\end{frame}

\end{document}
